\documentclass[10pt, letterpaper]{article}

% Packages:
\usepackage[
    ignoreheadfoot, % set margins without considering header and footer
    top=2 cm, % separation between body and page edge from the top
    bottom=2 cm, % separation between body and page edge from the bottom
    left=2 cm, % separation between body and page edge from the left
    right=2 cm, % separation between body and page edge from the right
    footskip=1.0 cm, % separation between body and footer
    % showframe % for debugging
]{geometry} % for adjusting page geometry
\usepackage[explicit]{titlesec} % for customizing section titles
\usepackage{tabularx} % for making tables with fixed width columns
\usepackage{array} % tabularx requires this
\usepackage[dvipsnames]{xcolor} % for coloring text
\definecolor{primaryColor}{RGB}{0, 79, 144} % define primary color
\usepackage{enumitem} % for customizing lists
\usepackage{fontawesome5} % for using icons
\usepackage{amsmath} % for math
\usepackage[
    pdftitle={David Ramirez CV},
    pdfauthor={David Ramirez Betancourth},
    pdfcreator={LaTeX with RenderCV},
    colorlinks=true,
    urlcolor=primaryColor
]{hyperref} % for links, metadata and bookmarks
\usepackage[pscoord]{eso-pic} % for floating text on the page
\usepackage{calc} % for calculating lengths
\usepackage{bookmark} % for bookmarks
\usepackage{lastpage} % for getting the total number of pages
\usepackage{changepage} % for one column entries (adjustwidth environment)
\usepackage{paracol} % for two and three column entries
\usepackage{ifthen} % for conditional statements
\usepackage{needspace} % for avoiding page break right after the section title
\usepackage{iftex} % check if engine is pdflatex, xetex or luatex

% Ensure generated PDF is machine readable/ATS parsable:
\ifPDFTeX
    \input{glyphtounicode}
    \pdfgentounicode=1
    \usepackage[T1]{fontenc}
    \usepackage[utf8]{inputenc}
    \usepackage{lmodern}
\fi

\usepackage[default, type1]{sourcesanspro}

% Some settings:
\AtBeginEnvironment{adjustwidth}{\partopsep0pt} % remove space before adjustwidth environment
\pagestyle{empty} % no header or footer
\setcounter{secnumdepth}{0} % no section numbering
\setlength{\parindent}{0pt} % no indentation
\setlength{\topskip}{0pt} % no top skip
\setlength{\columnsep}{0.15cm} % set column separation
\makeatletter
\let\ps@customFooterStyle\ps@plain % Copy the plain style to customFooterStyle
\patchcmd{\ps@customFooterStyle}{\thepage}{
    \color{gray}\textit{\small David Ramirez Betancourth - Page \thepage{} of \pageref*{LastPage}}
}{}{} % replace number by desired string
\makeatother
\pagestyle{customFooterStyle}

\titleformat{\section}{
    % avoid page break right after the section title
    \needspace{4\baselineskip}
    % make the font size of the section title large and color it with the primary color
    \Large\color{primaryColor}
}{
}{
}{
    % print bold title, give 0.15 cm space and draw a line of 0.8 pt thickness
    % from the end of the title to the end of the body
    \textbf{#1}\hspace{0.15cm}\titlerule[0.8pt]\hspace{-0.1cm}
}[] % section title formatting

\titlespacing{\section}{
    % left space:
    -1pt
}{
    % top space:
    0.3 cm
}{
    % bottom space:
    0.2 cm
} % section title spacing

\newenvironment{highlights}{
    \begin{itemize}[
        topsep=0.10 cm,
        parsep=0.10 cm,
        partopsep=0pt,
        itemsep=0pt,
        leftmargin=0.4 cm + 10pt
    ]
}{
    \end{itemize}
} % new environment for highlights

\newenvironment{highlightsforbulletentries}{
    \begin{itemize}[
        topsep=0.10 cm,
        parsep=0.10 cm,
        partopsep=0pt,
        itemsep=0pt,
        leftmargin=10pt
    ]
}{
    \end{itemize}
} % new environment for highlights for bullet entries

\newenvironment{onecolentry}{
    \begin{adjustwidth}{
        0.2 cm + 0.00001 cm
    }{
        0.2 cm + 0.00001 cm
    }
}{
    \end{adjustwidth}
} % new environment for one column entries

\newenvironment{twocolentry}[2][]{
    \onecolentry
    \def\secondColumn{#2}
    \setcolumnwidth{\fill, 4.5 cm}
    \begin{paracol}{2}
}{
    \switchcolumn \raggedleft \secondColumn
    \end{paracol}
    \endonecolentry
} % new environment for two column entries

\newenvironment{threecolentry}[3][]{
    \onecolentry
    \def\thirdColumn{#3}
    \setcolumnwidth{1 cm, \fill, 4.5 cm}
    \begin{paracol}{3}
    {\raggedright #2} \switchcolumn
}{
    \switchcolumn \raggedleft \thirdColumn
    \end{paracol}
    \endonecolentry
} % new environment for three column entries

\newenvironment{header}{
    \setlength{\topsep}{0pt}\par\kern\topsep\centering\color{primaryColor}\linespread{1.5}
}{
    \par\kern\topsep
} % new environment for the header

\newcommand{\placelastupdatedtext}{
  \AddToShipoutPictureFG*{
    \put(
        \LenToUnit{\paperwidth-2 cm-0.2 cm+0.05cm},
        \LenToUnit{\paperheight-1.0 cm}
    ){\vtop{{\null}\makebox[0pt][c]{
        \small\color{gray}\textit{Last updated in July 2026}\hspace{\widthof{Last updated in September 2024}}
    }}}%
  }%
}%

\let\hrefWithoutArrow\href
\renewcommand{\href}[2]{\hrefWithoutArrow{#1}{\ifthenelse{\equal{#2}{}}{ }{#2 }\raisebox{.15ex}{\footnotesize \faExternalLink*}}}

\begin{document}
    \newcommand{\AND}{\unskip
        \cleaders\copy\ANDbox\hskip\wd\ANDbox
        \ignorespaces
    }
    \newsavebox\ANDbox
    \sbox\ANDbox{}

    \placelastupdatedtext
    \begin{header}
        \fontsize{30 pt}{30 pt}
        \textbf{David Ramirez Betancourth}

        \vspace{0.3 cm}

        \normalsize
        \mbox{{\footnotesize\faMapMarker*}\hspace*{0.13cm}Manizales, Colombia}%
        \kern 0.25 cm%
        \AND%
        \kern 0.25 cm%
        \mbox{\hrefWithoutArrow{mailto:dramirezbe@unal.edu.co}{{\footnotesize\faEnvelope[regular]}\hspace*{0.13cm}dramirezbe@unal.edu.co}}%
        \kern 0.25 cm%
        \AND%
        \kern 0.25 cm%
        \mbox{\hrefWithoutArrow{tel:+57-310-680-13-31}{{\footnotesize\faPhone*}\hspace*{0.13cm}+57 3106801331}}%
        \kern 0.25 cm%
        \AND%
        \kern 0.25 cm%
        \mbox{\hrefWithoutArrow{https://linkedin.com/in/david-ramirez-betancourth-b5ba80279}{{\footnotesize\faLinkedinIn}\hspace*{0.13cm}david-ramirez-betancourth}}%
        \kern 0.25 cm%
        \AND%
        \kern 0.25 cm%
        \mbox{\hrefWithoutArrow{https://github.com/dramirezbe}{{\footnotesize\faGithub}\hspace*{0.13cm}dramirezbe}}%
    \end{header}

    \vspace{0.3 cm - 0.3 cm}

    \section{Professional Profile}

        \begin{onecolentry}
            Electronic Engineer specialized in Radio Frequency (RF) systems and Digital Signal Processing (DSP). Experience in hardware-software integration from low-level protocols (TCP, UDP, ZMQ) and Software Defined Radio (SDR) to cloud architectures with IoT sensors. Firmware development in C for real-time acquisition and processing, with autonomous calibration and Python orchestration. I complement my development workflow with AI agents for TDD/SDD automation, focusing on high-performance embedded architecture optimization.
        \end{onecolentry}

        \vspace{0.2 cm}

    \section{Education}

        \begin{threecolentry}{\textbf{MSc}}{
            Jan 2026 - Present
        }
            \textbf{National University of Colombia}, Early admission to M.Sc. in Industrial Automation
            \begin{highlights}
                \item \textbf{Research Areas:} Artificial Intelligence, Machine Learning applied to RF, Computational Structures, and Optimization.
            \end{highlights}
        \end{threecolentry}

        \begin{threecolentry}{\textbf{BS}}{
            Jul 2022 - Jun 2026
        }
            \textbf{National University of Colombia}, Electronic Engineering
            \begin{highlights}
                \item \textbf{Cumulative GPA:} 4.2/5.0
                \item \textbf{Focus Areas:} Digital Signal Processing (DSP), Software Defined Radio (SDR), Hardware Design, and IoT Systems.
            \end{highlights}
        \end{threecolentry}

    \section{Experience}

        \begin{twocolentry}{
            Manizales, Colombia

        Oct 2024 -- Present
        }
            \textbf{GCPDS Research Group}, Researcher and Developer
            \begin{highlights}
                \item C/Python firmware development for an IoT spectral monitoring sensor based on HackRF SDR on Raspberry Pi 5, with hybrid architecture (C for real-time IQ acquisition, Python for orchestration), autonomous PPM calibration, and WebRTC streaming.
                \item Backend platform implementation with REST API (FastAPI), TimescaleDB time-series database, and real-time spectral visualization dashboard.
                \item Mobile application development for deployment and validation of Image Segmentation models.
            \end{highlights}
        \end{twocolentry}

        \vspace{0.2 cm}

        \begin{twocolentry}{
            Manizales, Colombia

            2024
        }
            \textbf{National University of Colombia}, Undergraduate Fellow
            \begin{highlights}
                \item Academic teaching assistant for Signals and Systems and Machine Learning Theory, advising groups of 30+ students in theoretical and practical components.
                \item Design and instruction of hands-on Python workshops for signal analysis and processing, integrating tools such as NumPy, SciPy, and Matplotlib.
            \end{highlights}
        \end{twocolentry}

    \section{Projects}

        \begin{twocolentry}{
            \href{https://github.com/dramirezbe/SDR-SpectrumMonitoring-Sensor}{Repo: SDR-Sensor}
        }
            \textbf{IoT Spectrum Monitoring Sensor}
            \begin{highlights}
                \item Firmware development for a spectral monitoring system integrating Raspberry Pi 5 and HackRF SDR. Includes implementation of calibration protocols, precise NTP synchronization, and optimized drivers for SDR hardware.
                \item \textbf{Tools:} C, Python, HackRF, Raspberry Pi.
            \end{highlights}
        \end{twocolentry}

        \vspace{0.2 cm}

        \begin{twocolentry}{
            \href{https://github.com/dramirezbe/RPI-DashWake.git}{Repo: RPI-DashWake}
        }
            \textbf{RPI-DashWake (Smart Clock)}
            \begin{highlights}
                \item Design of a smart clock with integrated environmental monitoring (temperature, humidity, and air quality). Includes a web server for real-time metric visualization and programmable alarm management.
                \item \textbf{Technologies:} C, Python, Arduino, JavaScript.
            \end{highlights}
        \end{twocolentry}

        \vspace{0.2 cm}

        \begin{twocolentry}{
            \href{https://github.com/dramirezbe/SelfHostedContextCrafter}{Repo: ContextCrafter}
        }
            \textbf{SelfHostedContextCrafter (AI Agent Orchestrator)}
            \begin{highlights}
                \item Automated architectural team that orchestrates RAG-based expert agents, running fully on local hardware. It digests hardware datasheets, software documentation, and network protocols to generate structured context and automate development workflows.
                \item \textbf{Technologies:} Python, RAG, local LLMs, LaTeX.
            \end{highlights}
        \end{twocolentry}

    \section{Technical Achievements}

        \begin{onecolentry}
            \begin{highlightsforbulletentries}
                \item \textbf{Agnostic SDR architecture:} Architecture design for a cognitive sensor network, integrating WebRTC for real-time streaming, TimescaleDB for historical data, and high-performance WebGL visualization.
                \item \textbf{DSP and SDR calibration:} Development of iterative calibration algorithms (gradient descent) to correct clock offset (PPM) in SDR hardware, and implementation of custom DSP engines in C99 using Polyphase Filter Banks (PFB) and \textit{liquid-dsp}.
                \item \textbf{Telemetry and profiling on Pi 5:} Execution of hardware stress tests and power profiling (PD/PPS configurations) on Raspberry Pi 5, identifying and mitigating voltage drops under concurrent LTE and radiofrequency processing workloads.
                \item \textbf{AI automation (TDD/SDD):} Development workflow automation using AI agents with TDD and SDD methodologies, generating code, tests, and documentation predictively. Significant reduction in repetitive boilerplate and validation tasks.
                \item \textbf{Agent orchestrator (SelfHostedContextCrafter):} Development of a system that orchestrates RAG-based expert agents on local hardware, capable of digesting hardware datasheets and technical documentation to generate structured context and automate code review and deployment tasks.
            \end{highlightsforbulletentries}
        \end{onecolentry}

        \vspace{0.2 cm}

    \section{Technologies}

        \begin{onecolentry}
            \textbf{Languages:} C99, C++, Python, JavaScript, TypeScript, Verilog, SQL
        \end{onecolentry}

        \vspace{0.2 cm}

        \begin{onecolentry}
            \textbf{Frameworks and Platforms:} FastAPI, Django, React, Vite, Node.js, GNU Radio, SoapySDR, HackRF One, FPGAs (Gowin EDA)
        \end{onecolentry}

        \vspace{0.2 cm}

        \begin{onecolentry}
            \textbf{Infrastructure and Tools:} Linux (systemd, Bash), Docker, CMake, ZMQ, WebRTC, TimescaleDB, Git, Sphinx, Doxygen, NTP, Raspberry Pi, GPS/LTE
        \end{onecolentry}

    \section{Languages}

        \begin{onecolentry}
            \textbf{Spanish:} Native \quad|\quad \textbf{English:} Intermediate (B1)
        \end{onecolentry}

\end{document}
